% =============================================================================
% 4. STAGE SPECIFICATIONS
% Target: 4.5 pages (~2250 words)
% Source: cap4 §4.2–§4.5 (sec:stage1..sec:stage4)
% Status: SKELETON — keep only one operation + one property + one equation
%         per stage; move detailed tables and rework analysis to supplementary.
% =============================================================================

\section{Stage Specifications}
\label{sec:stages}

% =============================================================================
% 4.1 STAGE 1 — Event Log Construction (1 p)
% =============================================================================
\subsection{Stage 1: Event Log Construction}
\label{sec:stages-s1}

Stage~1 transforms heterogeneous repository data into a structured event
log. It is the composition
\begin{equation}
\label{eq:s1-decomposition}
\mathcal{S}_1 = \textsc{Format} \circ \textsc{Map} \circ
    \textsc{Correlate} \circ \textsc{Extract}.
\end{equation}

A raw event is a tuple
$e^{\text{raw}} = (ts, src, type, payload)$ with timestamp, source,
type and source-specific attribute map; the raw event pool is
$E^{\text{raw}} = \bigcup_{i,j} \textsc{Extract}_{src(D_i^{(j)})}(D_i^{(j)})$
with source-specific extractors for Git, issue tracker, CI/CD and
code review (full extraction/mapping tables in the online
supplementary).

\paragraph{Case correlation.}
The correlation function
$\kappa : E^{\text{raw}} \to \mathcal{C} \cup \{\bot\}$ is implemented
as a cascade of three strategies in decreasing specificity:
$\kappa_1$ explicit reference (issue keys in commit messages or PR
descriptions),
\begin{equation}
\label{eq:kappa1}
\kappa_1(e) = \begin{cases}
    \text{issue\_key}(e.\text{payload}) & \text{if extractable} \\
    \bot & \text{otherwise};
\end{cases}
\end{equation}
$\kappa_2$ structural link (platform-provided PR$\to$issue,
build$\to$commit, sub-task$\to$parent);
$\kappa_3$ heuristic (temporal proximity, author identity, file
overlap). Correlation quality is monitored by
\begin{align}
\text{Coverage} &= |E^{\text{corr}}| / |E^{\text{raw}}|
    \label{eq:corr-coverage} \\
\text{AmbiguityRate} &= |\{e : |\kappa(e)| > 1\}| / |E^{\text{raw}}|
    \label{eq:corr-ambiguity}.
\end{align}

\begin{definition}[Structured Event]
\label{def:structured-event}
A structured event is
$e = (\kappa(e), \mu(e), \tau(e), \rho(e), \omega(e))$ with case ID,
activity, timestamp, resource, and attribute map.
\end{definition}

After bot filtering and per-triple $(case, activity, timestamp)$
deduplication, the output event log is
\begin{equation}
\label{eq:final-log}
\mathcal{L} = \left\{ \sigma_c = \langle e_1, \ldots, e_{|\sigma_c|} \rangle
    \;\middle|\; c \in \mathcal{C}, \;
    \tau(e_k) \leq \tau(e_{k+1}) \; \forall k \right\}.
\end{equation}

% =============================================================================
% 4.2 STAGE 2 — Process Mining (1 p)
% =============================================================================
\subsection{Stage 2: Process Mining}
\label{sec:stages-s2}

Stage~2 produces a process model $M$, a metrics vector $\Phi$ and the
state space $S = S_T \cup S_A$ used by Stage~3. The Inductive
Miner~\cite{leemans2013discovering} is applied with noise threshold
$\theta_{\text{noise}} \in [0,1]$:
\begin{equation}
\label{eq:discovery}
M = \textsc{InductiveMiner}(\mathcal{L}, \theta_{\text{noise}}).
\end{equation}

\begin{property}[Soundness Guarantee]
\label{prop:soundness}
The Inductive Miner produces a sound process tree by
construction~\cite{leemans2013discovering}.
\end{property}

Alignment-based conformance produces
$\phi_{\text{fit}}, \phi_{\text{prec}}, \phi_{\text{gen}} \in [0,1]$
together with the per-case score
\begin{equation}
\label{eq:case-conformance}
\phi_{\text{conf}}(c) = 1 - \frac{\text{cost}(\text{align}(\sigma_c, M))}
    {\text{cost}(\text{worst\_align}(\sigma_c, M))},
\end{equation}
which feeds feature~5 of the ML component
(Section~\ref{sec:ml-component}). Variant entropy
$H_{\text{norm}}(V) = H(V) / \log_2 |V| \in [0,1]$ is computed as a
scale-independent complexity proxy. The state space is
$S = S_T \cup S_A$ with $S_A = \{\texttt{done}, \texttt{cancelled}\}$
and $S_T$ the non-terminal observed activities; the per-state mean
sojourn $\bar{d}_a$ feeds the expected lead time in
Equation~\ref{eq:expected-lead-time}. Per-activity variance, median
sojourn and per-transition waiting time are reported in the online
supplementary.

% =============================================================================
% 4.3 STAGE 3 — Stochastic Modeling (1.5 p)
% =============================================================================
\subsection{Stage 3: Stochastic Modeling}
\label{sec:stages-s3}

Stage~3 is the formal core of PATHCAST: it estimates a transition
matrix, an empirical sojourn-time distribution and the fundamental
quantities of an absorbing Markov chain. The decomposition is
\begin{equation}
\label{eq:s3-decomposition}
\mathcal{S}_3 = \textsc{Stationary} \circ \textsc{Sojourn} \circ
    \textsc{Absorb} \circ \textsc{Estimate} \circ \textsc{Count}.
\end{equation}

\paragraph{Maximum-likelihood estimation with optional smoothing.}
Let $C_{ij}$ be the count of transitions from state $i$ to state $j$
observed in $\mathcal{L}$. The smoothed MLE estimator is
\begin{equation}
\label{eq:mle-full}
\hat{P}_{ij} = \frac{C_{ij} + \alpha}
                    {\sum_{j'} C_{ij'} + \alpha |S|}, \qquad
\alpha \geq 0,
\end{equation}
with $\alpha = 0$ recovering the standard MLE. Symmetric smoothing
($\alpha > 0$) corresponds to a uniform Dirichlet prior on the row
simplex and stabilizes estimates for low-count rows.

\begin{property}[MLE Consistency]
\label{prop:mle-consistency}
For $\alpha = 0$, $\hat{P} \xrightarrow{a.s.} P^{*}$ as
$|\mathcal{L}| \to \infty$~\cite{Spedicato2017}.
\end{property}

\paragraph{Absorbing-chain decomposition.}
States are reordered so that transient precede absorbing:
\begin{equation}
\label{eq:absorbing-decomposition}
\hat{P} = \begin{pmatrix} Q & R \\ \mathbf{0} & I \end{pmatrix},
\qquad
N = (I-Q)^{-1}, \qquad
B = NR,
\end{equation}
yielding the expected number of visits matrix $N$ and the absorption
probability matrix $B$.

\begin{property}[Invertibility]
\label{prop:invertibility}
$(I-Q)$ is invertible iff every transient state can reach an absorbing
state~\cite{kemeny1976finite}.
\end{property}

\begin{property}[Absorption Certainty]
\label{prop:absorption-certainty}
$\sum_{a \in S_A} B_{ia} = 1$ for every $i \in S_T$.
\end{property}

Expected lead time in calendar time:
\begin{equation}
\label{eq:expected-lead-time}
\mathbb{E}[\text{LeadTime} \mid s_0 = i]
  = \sum_{j \in S_T} N_{ij} \cdot \bar{d}_j.
\end{equation}

\paragraph{Empirical sojourn-time distributions.}
For each $s \in S_T$, the observed sojourn set is
$\mathcal{T}_s = \{\tau(e_{k+1}) - \tau(e_k) \mid \mu(e_k) = s\}$
with empirical CDF
\begin{equation}
\label{eq:ecdf}
\hat{F}_s(t) = \tfrac{1}{|\mathcal{T}_s|}
    \sum_{\tau \in \mathcal{T}_s} \mathbf{1}[\tau \leq t].
\end{equation}

\begin{remark}[Non-Parametric Design Choice]
\label{rem:nonparametric}
Empirical resampling preserves the multi-modal, heavy-tailed and
zero-inflated patterns commonly observed in software process
sojourns, avoiding the misspecification incurred by lognormal,
Weibull or gamma fits~\cite{vacanti2015}.
\end{remark}

\paragraph{Confidence assessment.}
For row $i$ with $n_i = \sum_j C_{ij}$ observations,
\begin{equation}
\label{eq:se-transition}
\text{SE}(\hat{P}_{ij}) = \sqrt{\hat{P}_{ij}(1 - \hat{P}_{ij}) / n_i}.
\end{equation}
Rows with $n_i < n_{\min}$ (default 30) are flagged; the row-level
confidence grade is
\begin{equation}
\label{eq:confidence-grade}
\text{Grade}(P) = |\{i \in S_T : n_i \geq n_{\min}\}| / |S_T|,
\end{equation}
and $\text{Grade}(P) < 0.8$ is reported as a validity caveat. The
quasi-stationary distribution and rework-intensity decomposition are
reported in the online supplementary; full sensitivity to $\alpha$ is
analysed there.

% =============================================================================
% 4.4 STAGE 4 — Monte Carlo Simulation (1 p)
% =============================================================================
\subsection{Stage 4: Process-Aware Monte Carlo Simulation}
\label{sec:stages-s4}

Stage~4 generates the empirical forecast distribution by simulating
trajectories of the absorbing chain with sojourn times sampled from
$\hat{F}_s$:
\begin{equation}
\label{eq:forecast-distribution}
\hat{\mathcal{D}}
  = \{(T^{(m)}, o^{(m)}, \psi^{(m)})\}_{m=1}^{M_{\text{sim}}}.
\end{equation}

\begin{algorithm}[htbp]
\caption{Process-Aware Monte Carlo Simulation}
\label{alg:mc-full}
\begin{algorithmic}[1]
\REQUIRE $P$, $\hat{F}$, $\pi_0$ (or fixed $s_0$), $M_{\text{sim}}$,
    max steps $K_{\max}$
\ENSURE $\hat{\mathcal{D}}
    = \{(T^{(m)}, o^{(m)}, \psi^{(m)})\}_{m=1}^{M_{\text{sim}}}$
\STATE $\hat{\mathcal{D}} \gets \emptyset$
\FOR{$m = 1$ to $M_{\text{sim}}$}
    \STATE $s \gets s_0$ if fixed, else $s \gets$ sample from $\pi_0$
    \STATE $T \gets 0$; $\psi \gets [s]$; $k \gets 0$
    \WHILE{$s \notin S_A$ \AND $k < K_{\max}$}
        \STATE $\Delta t \gets$ sample from $\hat{F}_s$
        \STATE $T \gets T + \Delta t$
        \STATE $s' \gets \textsc{Categorical}(P[s, :])$
        \STATE $s \gets s'$; append $s$ to $\psi$; $k \gets k + 1$
    \ENDWHILE
    \IF{$s \notin S_A$}
        \STATE mark trajectory $m$ as \texttt{non-absorbed}
    \ENDIF
    \STATE $\hat{\mathcal{D}} \gets \hat{\mathcal{D}}
        \cup \{(T, s, \psi)\}$
\ENDFOR
\RETURN $\hat{\mathcal{D}}$
\end{algorithmic}
\end{algorithm}

The safety bound $K_{\max} = 1000$ prevents infinite loops in
degenerate cases; non-absorbed trajectories are excluded from
forecast aggregation. PATHCAST uses Mersenne Twister (MT19937) as the default PRNG and
records seeds per run for exact replication.

\paragraph{Forecast extraction.}
From $\hat{\mathcal{D}}$ the pipeline extracts the percentile forecasts
\begin{equation}
\label{eq:percentile}
\hat{Q}_\alpha = \inf\{t : \hat{F}_T(t) \geq \alpha\}
\quad (\alpha \in \{0.5, 0.85, 0.95\}),
\end{equation}
the completion probability
\begin{equation}
\label{eq:completion-prob}
\hat{P}(\text{Done})
  = |\{m : o^{(m)} = \texttt{done}\}| / M_{\text{sim}},
\end{equation}
and the tail risk
\begin{equation}
\label{eq:tail-risk}
\hat{P}(\text{late})
  = |\{m : T^{(m)} > \tau_{\text{deadline}}\}| / M_{\text{sim}}.
\end{equation}

\paragraph{Convergence.}

\begin{definition}[Convergence Criterion]
\label{def:convergence}
The simulation budget $M_{\text{sim}}$ satisfies the convergence
criterion at tolerance $\epsilon$ when
\begin{equation}
\label{eq:convergence-criterion}
\frac{z_{0.025} \cdot \hat{\sigma}_T / \sqrt{M_{\text{sim}}}}
     {\hat{Q}_{0.5}} < \epsilon.
\end{equation}
\end{definition}

\begin{property}[Convergence Rate]
\label{prop:convergence-rate}
The MC estimation error decreases as
$O(1/\sqrt{M_{\text{sim}}})$~\cite{robert2004monte}.
\end{property}

Adaptive doubling from $M_{\text{sim}} = 1{,}000$ to $100{,}000$
stops when Equation~\ref{eq:convergence-criterion} holds with
$\epsilon = 0.02$.
